\begin{document}
\headline{ACADEMIC CURRICULUM VITAE}{Computational scientist · Multiscale modeller}
\section{Research profile}
Applied mathematician and computational biologist working on mechanistic models across molecular, cellular and organismal scales. Research spans enzyme kinetics, nonlinear metabolic dynamics, stochastic immune-cell populations and plant–environment interactions. I connect model formulation to numerical simulation, parameter estimation, sensitivity, uncertainty and reproducible scientific software.
\section{Research experience}
\role{Researcher · Computational Cell Biology}{Oct 2026–present}{Institute for Computational Cell Biology (CCB), Heinrich Heine University Düsseldorf}
Research appointment resumed on 1 October 2026. Continuing work in computational biology and mechanistic modelling.
\role{Independent research \& scientific software}{Nov 2025–Sep 2026}{Düsseldorf, Germany}
Developed FokoLab and maintained research modelling projects; continued scientific writing and further training in AI engineering and research software.
\role{Postdoctoral researcher}{Nov 2022–Oct 2025}{CEPLAS / CCB, Heinrich Heine University Düsseldorf · Martin J. Lercher}
\begin{itemize}
\item Developed a thermo-hydraulic-biochemical modelling framework for photosynthetic adaptation across the C3–C4 continuum, linking resource allocation, heat exchange, water transport and environmental forcing.
\item Implemented simulation, parameter-exploration, optimisation and sensitivity workflows; worked with experimental collaborators and supervised student research.
\end{itemize}
\role{Marie Skłodowska-Curie doctoral researcher}{Aug 2019–Aug 2022}{PoLiMeR ITN / Institute of Quantitative and Theoretical Biology, HHU}
\begin{itemize}
\item Built kinetic ODE models of hepatic fatty-acid synthesis and metabolism; investigated equilibria, stability and conditions for metabolic switching.
\item Combined enzyme-kinetic literature, mathematical analysis and computational methods; contributed to peer-reviewed work on kinetic modelling and differentiable constraint-based models.
\end{itemize}
\role{Research fellow}{Nov 2017–Mar 2019}{African Institute for Mathematical Sciences · Cameroon and Rwanda}
Developed deterministic, delay and stochastic T-cell population models and inference workflows for division-tracking data; taught mathematical and computational subjects.
\newpage\continued{Research \& publications}
\section{Selected research and software}
\project{FokoLab · simulation, inference and scientific visualisation}{Browser-based scientific environment connecting custom models, dynamical systems, stochastic simulations, optimisation, sensitivity, fitting, statistics and mechanistic animation. Emphasis on recorded-state playback, explicit model scope and reproducible exports.\\\href{https://chilperic.github.io/chilperic_ode_solver/site/index.html}{Live platform} · \href{https://github.com/chilperic/chilperic_ode_solver}{Source repository}}
\project{Photosynthetic adaptation · mechanisms across scales}{Plant physiology framework joining biochemical assimilation, heat balance, hydraulics, soil resources and climate forcing. Student and collaborative work includes Jérémie Muller-Prokob and Yvonne Danisch as main collaborators, Antonio Rigueiro and Jonas Brass, under Martin J. Lercher at CCB.\\\href{https://github.com/chilperic/photosynthesis-climate-adaptation-model}{Research code}}
\project{Fatty-acid dynamics · kinetics and nonlinear behaviour}{Doctoral modelling of de novo synthesis and metabolic pools, including the structure of steady states and possible switching regimes. Scientific collaboration with Oliver Ebenhöh, Adélaïde Raguin and Barbara M. Bakker.\\\href{https://github.com/chilperic/fatty-acid-denovo-synthesis-model}{Synthesis model} · \href{https://github.com/chilperic/fatty-acid-metabolism-bistability}{Metabolic dynamics}}
\project{T-cell proliferation · populations and lineages}{Generation-structured models, Smith–Martin and Cyton formulations, and branching-process approaches to proliferation and loss. MSc research with Wilfred Ndifon and Gisèle Mophou; model-family provenance is kept distinct from later browser teaching implementations.\\\href{https://github.com/chilperic/tcell-proliferation-models}{Research code}}
\section{Publications}
\project{Preprint · 2025}{Hugo Dourado, \textbf{Chilperic Armel Foko Kuate}, Wolfram Liebermeister and Martin J. Lercher. \textit{Growth Mechanics and the emergence of metabolic oscillations in growing cells.} bioRxiv. \href{https://doi.org/10.1101/2025.06.24.661369}{\nolinkurl{doi:10.1101/2025.06.24.661369}}. Preprint.}
\project{Peer-reviewed review · 2023}{\textbf{Chilperic Armel Foko Kuate}, Oliver Ebenhöh, Barbara M. Bakker and Adélaïde Raguin. \textit{Kinetic data for modeling the dynamics of the enzymes involved in animal fatty acid synthesis.} Bioscience Reports 43(7), BSR20222496. \href{https://doi.org/10.1042/BSR20222496}{\nolinkurl{doi:10.1042/BSR20222496}}.}
\project{Peer-reviewed article · 2022}{St. Elmo Wilken, Mathieu Besançon, Miroslav Kratochvíl, \textbf{Chilperic Armel Foko Kuate}, Christophe Trefois, Wei Gu and Oliver Ebenhöh. \textit{Interrogating the effect of enzyme kinetics on metabolism using differentiable constraint-based models.} Metabolic Engineering 74, 72–82. \href{https://doi.org/10.1016/j.ymben.2022.09.002}{\nolinkurl{doi:10.1016/j.ymben.2022.09.002}}.}

\newpage\continued{Education, teaching \& development}
\section{Education}
\degree{PhD · Quantitative and Theoretical Biology}{2019–2023}{Heinrich Heine University Düsseldorf, Germany. Completed 7 March 2023.}
{\small Dissertation: \textit{Dynamic computational modeling of fatty acid de novo synthesis in the liver and qualitative analysis of fatty acid metabolism.} Supervision: Oliver Ebenhöh; co-supervision and collaboration: Adélaïde Raguin.}\par\vspace{4pt}
\degree{MSc · Mathematics and Fundamental Applications}{2017–2019}{Modelling specialisation · University of Yaoundé I, Cameroon.}
\degree{MSc · Mathematical Sciences}{2015–2016}{AIMS Ghana / University of Cape Coast, Ghana.}
\degree{BSc · Mathematics}{2009–2012}{University of Douala, Cameroon.}
\degree{Baccalauréat}{2009}{Collège La Confiance, Bafoussam, Cameroon.}

\section{Teaching, supervision and service}
\begin{itemize}
\item Supervised two MSc and two BSc research projects in computational biology, metabolism and plant physiology.
\item Delivered more than 200 hours of teaching and tutorials in mathematical modelling, differential equations, stochastic processes, algebra, scientific computing and Python across AIMS and HHU.
\item Early Career Researchers’ Representative, PoLiMeR ITN; earlier student representation at AIMS.
\end{itemize}
\section{Technical toolkit}
\textbf{Programming:} Python, Julia, R, MATLAB; JavaScript for scientific interfaces.\\
\textbf{Methods:} ODE/DAE models, branching processes, nonlinear dynamics, fitting, uncertainty, local/global sensitivity and constrained optimisation.\\
\textbf{Practice:} NumPy, SciPy, SALib, CasADi/IPOPT, CMA-ES; Git, Linux, Jupyter, LaTeX, testing and reproducible environments.
\section{Selected training and professional development}
\begin{tabularx}{\linewidth}{@{}p{24mm}X@{}}
2026 & AI-engineering and ML training in progress; scientific-software development; Bayer B4U mentoring programme.\\[5pt]
2023 & Optimal Control with CasADi, Leuven.\\[5pt]
2021 & Advanced AI and Data Science Workshop, HeiCAD.\\[5pt]
2020–2022 & PoLiMeR doctoral training programme.\\[5pt]
2018 & Epidemiological Modelling Clinic, AIMS South Africa / SACEMA.
\end{tabularx}
\section{Selected recognition}
Marie Skłodowska-Curie ITN fellowship, PoLiMeR; Post-AIMS Research Bursary; F.K.A. Allotey Meritorious Award, AIMS Ghana.
\section{Languages}
French: native · English: proficient · German: B1/B1+, with B2 study in progress.


\end{document}
